\chapter{Conclusions and Future Work}\label{ch:conclusions-and-future-work}

\section{Conclusions}\label{conclusions}

This thesis has documented the design and development state of Zephyr, a
four-actuator mobile robot whose Mecanum wheels are placed along a
common axis and whose narrow support geometry requires active pitch
stabilization. The work unifies mechanical configuration, electrical
architecture, embedded firmware, geometry-parameterized kinematics,
inverted-pendulum dynamics, control planning, and a traceable validation
program. The formulation intentionally avoids applying the familiar
rectangular four-corner Mecanum matrix to a geometry for which its lever
arms and roller ordering have not been verified.

Repository evidence confirms meaningful subsystem progress. Four AK40-10
actuators and 104 mm nominal Mecanum wheels have been selected; the
mechanical prototype is substantially assembled; IMU, receiver, display,
and CAN actuator interfaces operate in the ESP32 development firmware;
estimation and supervisor infrastructure exist; and a custom electrical
design is in development. A symbolic Jacobian, mobility-rank test,
nonlinear balance model, upright linearization, and constrained
allocation formulation provide the analytical basis for controller
design.

The central research question---whether the assembled collinear
mechanism can maintain balance while producing independently controlled
longitudinal, lateral, and yaw motion---cannot yet be answered
quantitatively. The exact wheel sequence, roller signs, wheel
coordinates, mass properties, actuator configuration, final control
gains, timing distribution, and synchronized test data remain required.
No claim of stable balance, holonomic mobility, disturbance tolerance,
tracking accuracy, power margin, or thermal margin is made in their
absence.

\par\medskip\noindent\singlespacing
\phantomsection\addcontentsline{lot}{table}{\protect\numberline{12.1}Objective completion at the time of this draft}
\textbf{Table 12.1. Objective completion at the time of this draft}\par

\begin{longtable}[]{@{}
  >{\raggedright\arraybackslash}p{(\columnwidth - 6\tabcolsep) * \real{0.2180}}
  >{\raggedright\arraybackslash}p{(\columnwidth - 6\tabcolsep) * \real{0.3083}}
  >{\raggedright\arraybackslash}p{(\columnwidth - 6\tabcolsep) * \real{0.1729}}
  >{\raggedright\arraybackslash}p{(\columnwidth - 6\tabcolsep) * \real{0.3008}}@{}}
\toprule\noalign{}
\begin{minipage}[b]{\linewidth}\raggedright
\textbf{Objective}
\end{minipage} & \begin{minipage}[b]{\linewidth}\raggedright
\textbf{Evidence completed}
\end{minipage} & \begin{minipage}[b]{\linewidth}\raggedright
\textbf{Status}
\end{minipage} & \begin{minipage}[b]{\linewidth}\raggedright
\textbf{Evidence still required}
\end{minipage} \\
\midrule\noalign{}
\endhead
\bottomrule\noalign{}
\endlastfoot
Document actual architecture & Repository-based subsystem and interface
record & Partially achieved & Final as-built survey and branch-complete
archive \\
Establish collinear kinematics & Symbolic geometry\slash{}sign model and rank
test & Analytically achieved & Measured y\_i, s\_i, radius, and
experimental validation \\
Model balance dynamics & Nonlinear and linear symbolic models &
Analytically achieved & Measured mass properties and identified
losses \\
Implement embedded infrastructure & IMU, CRSF, display, CAN, feedback,
supervisor paths & Substantially achieved & Release build, timing,
four-node configuration, logging \\
Implement stabilizing controller & Architecture and safety constraints
specified & Not yet demonstrated & Code, gains, constrained test logs \\
Demonstrate omnidirectional motion while balancing & Test method
specified & Not yet demonstrated & External pose and synchronized
vehicle data \\
Verify electrical and safety design & Components and design rules
documented & Partially achieved & Final schematic\slash{}layout, fabrication,
bench\slash{}thermal\slash{}fault tests \\
\end{longtable}

\section{Mechanical Future Work}\label{mechanical-future-work}

The highest-priority mechanical task is an as-built survey. It should
record wheel-center coordinates, wheel order, roller handedness,
effective loaded radius, frame dimensions, fasteners, shafts, bearings,
couplings, actuator mount stiffness, component masses, complete vehicle
mass, center of mass, and principal inertias. CAD should then be
reconciled with the assembled prototype. The design should provide
repeatable wheel alignment, protected cable routing, accessible
fasteners, a battery restraint, lifting points, and a fall structure
that protects hands, electronics, and rollers without interfering with
normal motion. Structural calculations and safety factors should be
added only after materials, loads, and interfaces are confirmed.

\section{Electrical Future Work}\label{electrical-future-work}

The schematic must be frozen only after the installed battery,
protection approach, connectors, fuse ratings, regulator range, CAN
transceiver, termination, and actuator peak demands are verified. The
PCB should be completed with a documented copper weight, stack-up,
branch-current assumptions, thermal paths, return-current control,
decoupling, transient protection, and design-rule review. Low-voltage
bring-up should precede battery testing, and the fabricated assembly
should undergo short-circuit inspection, rail-load testing, thermal
testing, CAN signal-integrity inspection, undervoltage response, and
fault-injection tests.

\section{Firmware and Control Future Work}\label{firmware-and-control-future-work}

Firmware priorities are to freeze a reproducible ESP32 baseline, verify
all four CAN node IDs and directions, measure task timing and feedback
age, implement structured logging, and complete the command-to-actuator
allocation. The controller should first be validated in restrained tests
with explicit saturation, rate limits, anti-windup, stale-data
rejection, tilt limits, and receiver failsafe. LQR or another
state-feedback law may be used for the inner balance loop if the
measured model is controllable and the estimator supplies reliable
states \cite{anderson1990optimal}; outer planar loops should be added only after balance
torque reserve is quantified. The STM32H723 migration should remain a
separate milestone until behavior is equivalent on the present platform.

\section{Experimental Future Work}\label{experimental-future-work}

The proposed experiment sequence is: as-built metrology; power and CAN
validation; IMU and actuator calibration; restrained balance; free
balance; repeatable disturbance recovery; independent longitudinal,
lateral, and yaw commands; combined trajectories; and endurance and
fault injection. Each stage should have approved quantitative
requirements and stop criteria before execution. External pose
measurement, independent electrical instrumentation, synchronized
clocks, repeated trials, uncertainty reporting, and retention of
failures are needed for a defensible model-to-hardware comparison.

When these steps are complete, the final thesis should replace the
provisional abstract, pending result language, and ready-to-fill tables
with measured values. The strongest potential contribution is not simply
a four-wheel robot, but a validated account of the mobility and control
consequences of placing four Mecanum wheels on one line while
simultaneously stabilizing an inverted body. Whether that potential is
realized remains an empirical question to be answered by the documented
test program.
